% HD-F1-I04. Root proof; independent review pending.
\section[HD-F1-I04: uniform reflection-rich full-group tuples]{HD-F1-I04: explicit full-group tuples throughout the previously proved reflection-rich rank range}
\label{sec:f1-reflection-rich-full-rank}

Retain the actual monomial tower, faithful primitive character,
specified polarized corner and balanced signature$(n,n)$ of
HD-I01--HD-I21. Write$g=g(Y)$,$b$ for the total branch count,
$\nu$ for the nonreflection count and$A=b-\nu$ for the
reflection count. HD-I16, under the exact numerical and morphism comparison below, already proves that every such component in the range
\[
                 A\ge 2n+2g-1                         \tag{F1I04.1}
\]
has, at every smooth Hurwitz point,
\[
 \operatorname{rank}d\mathcal F=2n+g-1,
 \qquad \ker T_h=0,
 \qquad \operatorname{rank}d\mathrm{Per}
                        =\operatorname{rank}\mu.       \tag{F1I04.2}
\]
Here$\mathcal F$ is the original cyclic-data forgetful map,
$T_h$ its actual quadratic trace and$\mu$ the original
primitive eigendifferential multiplication. For every retained
genus$g\le n-2$, these equalities give the complete all-point
period rank$2n+g-1$. For the upper three genera, the same prior criterion identifies
the period rank with that actual multiplication rank.

The rational class is nonempty throughout a full interval:
for every$n\ge5$, every permitted
$m\ge\lceil(n^2-2n+3)/2\rceil$, and every integer
\[
           2n-1\le A\le 2m+2n-8,                    \tag{F1I04.3}
\]
there is an actual full$G(3,1,m)$ component with$g=0$,
$\nu=4$, reflection count$A$, and period rank$2n-1$
at every smooth point. Its forgetful tangent kernel has
dimension$A-2n+2$. In particular, for$n=5,m=9$, every
$A=9,\ldots,20$ occurs and the original period rank is nine.

\paragraph{1. Exact identification with the previously proved range.}
The rank assertion in (F1I04.2) is the already proved HD-I16
criterion, expressed in reflection counts. It is retained as
an input here, not counted as a new rank theorem.
Indeed, with the original moving-rank excess of I16,
\[
 N=2m+2n-b,\qquad A=b-\nu,\qquad
 N+\nu=2m+2n-A.
 \tag{F1I04.4}
\]
Consequently its inequality $N+\nu\le2m-2g+1$ is exactly
(F1I04.1). I16 paragraph 3 proves that the original
forgetful map is onto the cyclic deformation space at
every smooth point in this range. I20 identifies the dual
of that actual differential with the canonical-sheaf
quadratic trace, whose local expression is
\[
 f(z)dz^2\longmapsto
       \sum_{h(z)=w}\frac{f(z)}{h'(z)^2}\,dw^2.
 \tag{F1I04.5}
\]
Thus the same differential gives $\ker T_h=0$.
The I21 pushout of the prescribed unit sequence identifies
its tangent kernel with $H^0(L\otimes\mathcal W_Q)$,
where $L=\mathcal O(2-\nu)$ and
$\mathcal W_Q=h_*\mathcal O_Y(Q)/\mathcal O$.
These are the actual maps of the given tower.

The original period codifferential is $T_h\mu$;
its rank therefore equals that of the actual primitive
multiplication $\mu$. The I02 comparison is the restriction
of the given $q^*$, with $q_*q^*=[|\ker\chi|]$ and its
specified theta pullback multiplier. It carries these
ranks to the given polarized Weil-period map. I12 makes
$\mu$ onto at every point for $g\le n-2$. In particular,
on a rational source it is ordinary multiplication of
sections of the actual two line bundles $\mathcal O(n-1)$.
The resulting rank $2n+g-1$, and its upper-genus
multiplication-rank form, were already consequences of
I16 and I20--I21. The additional assertion proved below
is the uniform existence of the explicit full-group tuples
in (F1I04.3).

\paragraph{2. Explicit full-group components throughout the interval.}
Use colors in$\mathbf Z/3$, and let$R_i(a)$ have the single
diagonal entry$\zeta_3^a$ at$i$. Put$k=m-3$ and define the
ordinary permutation matrices
\[
 C=(1\,2\,\cdots\,k),\qquad
 D=(k\,k+1\,k+2\,k+3),
\]
and actual diagonal matrices
\[
 U=R_{k+1}(1)R_{k+2}(1)R_{k+3}(1),\qquad
 V=R_k(1)\prod_{i=1}^{n-3}R_i(1).                    \tag{F1I04.11}
\]
The permitted bound gives$m\ge n+2$ and$k\ge6$, so the indices$1,\ldots,n-3$ are fixed by$D$;
the indices$k+1,k+2,k+3$ are fixed by$C$.
Start with the word
\[
       C,\ U,\ C^{-1},\ U^{-1},\
       D,\ V,\ D^{-1},\ DV^{-1}D^{-1},                \tag{F1I04.12}
\]
where each displayed diagonal matrix is expanded into its
individual nonidentity$R_i$ factors, in the written order.
Those factors commute within their indicated block.
The product is one: $C$ commutes with$U$, and the last
four blocks multiply to
$DVD^{-1}DV^{-1}D^{-1}=1$. There are four nonreflection
entries, namely$C,C^{-1},D,D^{-1}$, and$2n+2$ diagonal
reflections. These are actual matrices, not just conjugacy
classes with an unspecified product.

We prove full generation directly. The commutator of$C$
and$D$ is$(k\,1\,k+1)$. Its conjugates by powers of$C$
give all consecutive triples$(k+1\,i\,i+1)$ along the
$k$-cycle. These generate$A_{\{1,\ldots,k+1\}}$.
Indeed two triples sharing two letters generate the
alternating group on their four-letter union by conjugation;
then transitivity on ordered pairs of the old alternating
group supplies every triple with the next letter. Induction
along the cycle proves the assertion. Conjugating the
displayed commutator once by$D$ gives
$(k+1\,1\,k+2)$, and twice gives$(k+2\,1\,k+3)$.
The same two-old-letter argument gives$A_m$. Since$D$
is an odd permutation, $C,D$ generate$S_m$. One of the
individual diagonal reflections has color one; conjugating
it by the actual permutation matrices gives every$R_i(1)$.
They generate the full diagonal group. Thus(F1I04.12)
generates precisely$G(3,1,m)$.

The four permutation indices sum to
\[
               2(k-1)+2\cdot3=2m-2.                 \tag{F1I04.13}
\]
The diagonal reflections have permutation index zero, so
Riemann--Hurwitz makes the actual degree-$m$ quotient
rational. Exactly$n+1$ of its colored points have color
one and$n+1$ have color two. In either primitive eigenspace
the fractional-color sum is$n+1$; hence Chevalley--Weil
gives$-1+(n+1)=n$. Both multiplicities are the specified$n$.

To obtain$A=2n-1,2n,2n+1$, combine respectively three,
two or one of the first$U$ factors with the immediately
preceding$C$. The resulting entry remains a nonreflection,
since its uncolored$k$-cycle has moving rank$k-1\ge5$.
All other entries are unchanged, and the total product
is unchanged. The tuple still contains$C^{-1}$ and$D$,
so it contains the original$S_m$, as well as a color-one
diagonal reflection from$V$. Full generation persists.
Each combined factor becomes a colored fixed cycle in
that nonreflection fibre. The permutation ramification
and both fractional-color sums are exactly unchanged.

For every larger value$A\ge2n+2$, begin with the fully
expanded(F1I04.12). On any of its four ordinary cycles,
use the literal factorization
\[
 (a_1\,\cdots\,a_\ell)
   =(a_1\,\cdots\,a_{\ell-1})(a_{\ell-1}\,a_\ell)
 \qquad(\ell\ge4).                                   \tag{F1I04.14}
\]
Each use adds one uncolored transposition reflection and
leaves a nonreflection cycle of length at least three.
Apply it successively to the two length-$k$ cycles and
the two length-four cycles. Their total capacity is
$2(k-3)+2=2m-10$. Every integer between zero and this
capacity is attained by stopping the sequence after that
many uses. Thus every$A$ from$2n+2$ through
$2n+2+(2m-10)=2m+2n-8$ occurs.

Multiplying each new factor block recovers the original
entry. Therefore the generated group still contains the
old full group, and every new factor belongs to that same
$G(3,1,m)$. Formula(F1I04.14) preserves permutation index,
because$\ell-1=(\ell-2)+1$. All new cycles and reflection
fibres are uncolored, so the colored divisor and the two
fractional-color sums remain unchanged. There are exactly
four nonreflection values throughout.

Riemann existence applied to each explicit product-one
generating tuple gives a connected Galois cover and its
actual monomial quotient tower. Its branch values vary
in the ordinary smooth configuration space. The resulting
Hurwitz components are therefore nonempty with the asserted
invariants. The already proved all-point rank formula
applies to each such component, establishing(F1I04.3).

\paragraph{3. Published rational comparison through the original map.}
Khashayar Filom, \emph{On the Rank of the Thurston Pullback
Map}, IMRN2023, no.~24, 21838--21873,
doi:10.1093/imrn/rnac305, Theorem1.1, already gives the
rational rank criterion after the following precise
comparison. Mark the actual colored points by$R$ and all
target branch values by$B_{\rm all}$. The original$h$
is an admissible marked sphere cover in that theorem.
The local marked Hurwitz map to its branch configuration
is etale; composing its source-marking map with the local
Teichmuller marking is the theorem's pullback map. Its
cotangent operator is exactly the quadratic trace$T_h$
in paragraph1. These are the same maps used in I01--I03.

A diagonal reflection gives a regular value of$h$ with
exactly one marked source point. It contributes to Filom's
$\ell_1$. A transposition reflection has one simple critical
point, no other critical point in that fibre, and no colored
point there; it contributes to$\ell_2$. Hence
$\ell_1+\ell_2\ge A$. Since$|R|=2n+2$ for$g=0$, the
published lower bound is
\[
       \operatorname{rank}d\mathcal F
                   \ge\min(A,2n-1).                  \tag{F1I04.15}
\]
For$A\ge2n-1$ this is the full rank. The previously proved I16 cohomological argument supplies
the stated higher-genus range through the actual quotient maps. No
priority is claimed for the rational criterion.

The complete supplied proof of Theorem1.1, including
Lemma3.1, was read in the primary author version
arXiv2202.00768v2,14October2022, printed pages8--11.
It detects the possible poles of a kernel quadratic
differential and its required zeros at the indicated
simple critical points. In Lemma3.1 the two inverse
branches have local parameters$\pm\sqrt t$; their summed
analytic coefficient is even, giving the displayed
nonzero$C/t+O(1)$ term. The primary journal metadata was
separately checked. The retained higher-genus I16 proof uses the
actual Serre-dual line bundle and the checked I20--I21
morphisms; it imports no rational
Teichmuller assertion into a positive-genus source.

\paragraph{Framework and remaining scope.}
The finite covers and their original quotient maps are
unchanged. The retained I16 jet sequence and I21 unit pushout,
followed by canonical trace and the actual
polarized$q^*$, give the F01 variational map and its
specified tangent rank. The family of quotient maps and
marked divisors is the same Hilbert/Chow realization in
F02. The$q_*q^*$ identity supplies its F03 Poincare
biextension comparison with the stated multiplier.
Those retained coherent exact sequences give their actual
$K_0$ relations; the retained Adams and cycle-class maps
of F04 apply to those same classes. No detector cycle or
division constant is added by this component calculation.

The uniform tuples establish the additional nonemptiness claim
in (F1I04.3). The reflection-rich rank criterion itself was
already settled by I16. This result does not count that
criterion again as a narrowing of T4-INDIVIDUAL. Below that
range, the exact I21 bundle and multiplication-intersection
gate retains its previous scope.
